\chapter{前沿专题：人工智能、大模型与 Web3}

第 13 章的结尾把问题引向此处：人工智能的经济影响、大模型时代的产业组织、Web3 的治理挑战，这些命题的检验依赖计量工具，而它们的分析则依赖前 12 章的理论。本章的工作就是把工具搬上前沿：AI 的就业效应用第 9 章的任务模型分析，大模型的产业组织用第 3 章与第 4 章的规模经济框架分析，Web3 的代币经济用第 10 章的货币理论分析。本章的价值不在名词的新颖，而在分析的可迁移性：面对任何新技术叙事，先识别其经济结构，再调用已有工具，最后落到证据与政策，这是论述题与复试作答的通用路径。

\begin{examguide}
\textbf{本章考情}：本章为前沿素材章（\important{专硕·核心 + 学硕·热点}），为论述题与复试提供最前沿的分析工具：AI 的就业与增长效应、"人工智能+"政策、大模型产业组织、Web3 与代币经济是复试问答的高频话题。本章的价值在于用前 13 章的工具分析新现象。

\textbf{核心考点}：（1）AI 作为通用目的技术的证据与争论；（2）自动化边界的认知任务扩展与任务分析单位；（3）算法合谋的强化学习机制；（4）算法歧视的传导与治理；（5）大模型的规模报酬与训练成本；（6）开源与闭源竞争与 API 经济；（7）代币经济的货币数量方程与质押激励；（8）中国数字经济政策体系（2522、数据要素×、人工智能+）；（9）数字经济与中国式现代化。

\textbf{本章主线}：AI 的经济影响（14.1）→ 大模型产业组织（14.2）→ Web3 与代币经济（14.3）→ 政策与展望（14.4）。
\end{examguide}

\section{人工智能的经济学}

\subsection{AI 作为通用目的技术}

人工智能在经济学中的定位问题，直接决定其影响的分析框架：如果 AI 只是又一次普通的技术进步，那么第 8 章的增长核算与第 9 章的任务模型已经足够刻画它的效应；如果 AI 是与电力同级的通用目的技术，那么它的影响会经由互补投资、组织变革与技能结构调整的长链条逐步展开，短期证据必然低估长期影响。判断 AI 的性质因此是本章的第一个问题。

\begin{definition}[通用目的技术（GPT）]
\textbf{通用目的技术}（general purpose technology，GPT）指同时满足三个特征的技术：\textbf{广泛渗透性}，能应用于众多下游部门；\textbf{持续改进性}，技术本身长期持续进步；\textbf{创新互补性}，能催生下游部门的创新浪潮。
\end{definition}

按这三条标准逐项对照，AI 的证据分布并不均匀。第一条与第二条的正面证据已经充分：生成式 AI 在两年内渗透到软件开发、教育、医疗、金融、内容创作等数十个行业，模型能力以月为单位持续迭代；第三条仍在形成之中：AI 催生的下游创新（AI 原生应用、行业智能体）刚刚起步，能否构成一波类似电气化时代的创新浪潮尚无定论。反对论据集中在生产率一侧：宏观生产率数据尚未出现 AI 驱动的明显抬升，第 8 章讨论的索洛悖论在 AI 时代重现。支持一方对此的标准回应是时滞：电力在 1890 年代就完成关键技术突破，但制造业生产率的大幅提升要到 1920 年代才显现，因为工厂布局、组织流程与技能体系的改造需要数十年。AI 的投资与生产率红利之间同样存在互补投资时滞，当前证据不足不等于效应不存在。\mn{AI 的 GPT 判定（复试高频）：按三标准逐条论证：渗透性（已显现）、持续改进（已显现）、创新互补性（进行中）。争论焦点不在前两条，而在生产率红利的时滞长度。}从方法论上，这一争论与第 8 章通用目的技术时滞模型完全同构：判断 AI 是否"真正有用"，需要等待互补资本（组织变革、技能更新、流程再造）的积累完成，而互补资本的投资决策本身依赖对 AI 前景的预期，预期自我实现的逻辑（第 3 章）在技术扩散层面再次出现。

渗透性的直接证据在中国尤为显著：生成式人工智能用户规模已达约 6 亿（中国互联网络信息中心，2025），普及率超过四成，这一扩散速度远超个人计算机与互联网的早期阶段。也正因如此，2024 年政府工作报告将"人工智能+"行动与"互联网+"并置，其政策逻辑正是把 AI 当作通用目的技术来部署：不限定应用场景，而是让所有产业部门成为 AI 创新的下游（14.4 节展开）。AI 的 GPT 判定由此不仅是学术争论，更是一国技术政策的前提假设。

\subsection{AI 与劳动市场：任务模型的扩展}

第 9 章的任务模型把自动化刻画为边界 $\bar\theta$ 的外移，并留下"自动化与新任务创造竞赛"的核心命题，其末尾已经指出任务模型的要义在 AI 处变得锋利。AI 对这一框架的冲击不改变模型的结构，却改变两个关键参数的取值：边界的性质与边界的移动速度。本节把这两处变化展开为机制细节。

变化之一在边界的性质。历史上自动化沿"体力任务到程序任务"的顺序渐进推进，被自动化的任务集中在制造业装配、文书处理等可编码区间；生成式 AI 使认知任务首次进入自动化范围：文本撰写、代码生成、图像设计、初级法律文书，这些高技能劳动者安身立命的任务成为新的替代对象。边界 $\bar\theta$ 的外移不再蚕食低技能区间，而是直接切入高技能区间，替代效应的分布方向发生了逆转。变化之二在边界的移动速度。AI 能力的改进以月计（模型代际迭代），而劳动者的再技能化以年计，自动化边界的外移速度系统性快于技能供给的调整速度：

\[
\frac{d\bar\theta}{dt} \gg \frac{d(\text{技能供给})}{dt}.
\]

第 9 章的分析在此给出一个直接推论：均衡调整的传统机制是劳动者向更复杂的任务迁移，这一机制成立的前提是新任务的创造速度快于自动化边界的外移。当边界外移骤然加速而新任务创造的速度不变时，调整期内的摩擦性失业上升，就业冲击集中于那些任务构成中认知任务占比高的职业。这一推论可以形式化。

\begin{derivation}{AI 下自动化边界的外移与就业}{}
\textbf{第一步，设定任务结构。}沿用第 9 章的任务模型，把总任务测度标准化为 1，任务按复杂度 $x$ 排列，任务 $x$ 处资本的相对生产效率为 $\phi(x)$，且 $\phi'(x) > 0$：任务越简单（$x$ 小），资本越擅长。设每单位资本的使用成本为 $r$，工资为 $w$，任务 $x$ 由劳动承担的条件是 $w \le r\phi(x)$，解得临界任务
\[
x^* = \phi^{-1}\!\left(\frac{w}{r}\right).
\]
区间 $[0, x^*]$ 为自动化区间（由资本完成），$(x^*, 1]$ 由劳动完成，$x^*$ 即自动化边界。\textbf{第二步，引入 AI 的外移。}AI 通过改变 $\phi(x)$ 的斜率使边界加速外移，对时间求导并把边界速度与工资成本变化分离：
\[
\frac{dx^*}{dt} = \frac{1}{\phi'(x^*)}\left[\frac{\dot w}{w} - \frac{\dot r}{r}\right]\frac{w}{r}.
\]
AI 降低资本使用成本（$\dot r/r < 0$）且提升资本效率，括号项扩大，边界速度与 AI 能力改进同向。\textbf{第三步，就业冲击条件。}记新任务创造速度为 $g_N$，劳动侧任务区间的变化率为
\[
\frac{d(1-x^*)}{dt} = g_N - \frac{dx^*}{dt}.
\]
结果：\important{当自动化边界的外移速度超过新任务的创造速度时，劳动侧任务区间收缩，调整期内失业上升}；AI 对认知任务的切入使 $dx^*/dt$ 骤然变大，而新任务的识别与学习以年计（$g_N$ 几乎不变），于是 $g_N - dx^*/dt < 0$，就业冲击集中于认知任务占比高的职业。这就是第 9 章"自动化与新任务创造竞赛"在 AI 时代的形式化版本。
\end{derivation}

\begin{misconception}
"AI 将整体替代劳动"是常见误区。任务模型给出的正确分析单位是\textbf{任务}而非\textbf{职业}：职业是任务束，AI 替代部分任务的同时可能强化同一职业的剩余任务。以医生为例，AI 读片的精度超过初级影像科医生，但读片任务的自动化释放了医生的时间，使其专注于诊断沟通与治疗决策这些 AI 尚不能完成的任务。职业层面的净效应取决于任务束内替代任务与互补任务的相对权重，不能从"某项任务被替代"跳到"整个职业消失"。
\end{misconception}

任务层面的替代与职业层面的净效应分离，已经获得大样本证据的支持。Brynjolfsson 等（2023）在一家财富 500 强企业的 5000 余名客服人员中引入生成式 AI 助手，发现每小时解决的问题数平均提升约 14\%，其中低经验员工提升约 35\%，高经验员工几乎无提升：AI 把优秀员工的隐性知识扩散给了新员工，替代的是经验差距而非劳动者本身。中国劳动力市场的证据同样指向任务视角：2024 年人力资源社会保障部发布的新职业目录首次纳入生成式人工智能相关职业，"机器换人"与"新职业涌现"在同一个劳动力市场同时发生，这正是自动化与新任务创造竞赛的当代样本。\mn{AI 就业论述模板：证据先行（Brynjolfsson 等，2023，客服效率提升约 14\%、低经验员工约 35\%），结论落在"看任务、不看职业"，替代与互补并存。}\important{AI 就业分析的第一方法论原则：分析单位必须是任务而非职业}，替代效应与互补效应在任务层面并存，职业层面的净效应是二者加总的结果。

\subsection{算法合谋的强化学习机制}

AI 的另一个经济角色发生在市场内部：当定价决策从人转移到算法，合谋的达成条件随之改变。第 6 章用惩罚时滞模型说明算法如何稳定合谋：人工定价的报复时滞以周计，算法定价的报复时滞以秒计，超竞争价格的可持续性条件从可能违反变为恒成立。本节从学习机制的角度推进这一问题：定价算法甚至不需要被"教"会合谋，它可以在反复交互中自己学会合谋。

第 6 章的分析预设了合谋价格已经存在，算法只是让惩罚变得即时。\textbf{强化学习}（reinforcement learning）定价算法的出现则取消了这一预设。强化学习算法不需要预先编程的价格规则，它通过"试错"学习：每次报价后观察利润反馈，利润上升的策略被强化，利润下降的策略被淘汰。Calvano 等（2020）在受控实验中让两个 \textbf{Q-learning} 定价智能体进行双寡头价格竞争，结果两个算法在没有沟通、没有预设合谋指令的条件下，自发收敛到显著高于纳什均衡的超竞争价格，且对任何降价偏离实施价格战惩罚，惩罚结束后又自动回归超竞争价格。\important{算法合谋不需要协议，甚至不需要教唆：合谋是算法交互学习的均衡结果}，这是第 6 章四类形态中"自主机器"类从理论可能变为实验事实的标志。

从实验中可以归纳出强化学习定价算法收敛到合谋的三个条件。其一，\textbf{市场透明}：对手的价格必须实时可观测，学习信号才能即时更新；电商平台的价格抓取与比价网站恰好提供了这种透明度。其二，\textbf{交互频率高}：定价调整越频繁，偏离被发现和报复得越快，第 6 章时滞条件中的 $\Delta t$ 越趋近于零。其三，\textbf{学习能力足够}：算法必须有足够的探索空间去发现"提价有利、降价招致报复"的规律。三个条件在数字市场同时满足，合谋的达成因此不需要任何人类意图。三个条件各有一个形式化对应。

\begin{derivation}{算法定价与合谋的可持续性}{}
\textbf{第一步，重复博弈基准。}设合谋价格下每期利润为 $\pi^m$，单方面降价的偏离利润为 $\pi^d$（$\pi^d > \pi^m$），纳什竞争利润为 $\pi^n$（$\pi^n < \pi^m$）。触发策略下，合作的贴现值与偏离的贴现值分别为
\[
V^c = \frac{\pi^m}{1-\delta}, \qquad V^d = \pi^d + \frac{\delta\,\pi^n}{1-\delta},
\]
其中 $\delta$ 为下一期收益的贴现因子，$\delta$ 越大说明未来利润对当前决策的权重越高。合谋可持续的条件是 $V^c \ge V^d$，解得
\[
\delta \ge \delta^* = \frac{\pi^d - \pi^m}{\pi^d - \pi^n}.
\]
偏离的当期利得越大（$\pi^d$ 高）、惩罚后的竞争利润越接近偏离利润（$\pi^n$ 接近 $\pi^d$），所需的 $\delta$ 越高。\textbf{第二步，报复时滞进入贴现因子。}设价格调整的时间间隔为 $\Delta t$，连续复利利率为 $r$，则 $\delta = e^{-r\Delta t}$。人工定价下 $\Delta t$ 以周计，$\delta$ 显著小于 1；算法定价把 $\Delta t$ 压缩到秒级，$\delta \to 1$，而 $\delta^* \le 1$ 恒成立，合谋条件在极限情形下恒被满足：
\[
\Delta t \to 0 \quad\Longrightarrow\quad \delta \to 1 \ge \delta^*.
\]
\textbf{第三步，学习机制的重述。}Q-learning 定价算法不预设 $\pi^m$ 的位置，而是以 $\alpha$ 为学习率更新行动价值：$Q(p) \leftarrow Q(p) + \alpha[\pi(p) + \delta\max_{p'}Q(p') - Q(p)]$。在偏离会触发降价惩罚的环境中，$Q(p^m) > Q(p^d)$ 成立，算法逐步把合谋价识别为价值最高的行动。合谋价被发现需要足够高的 $\alpha$（学习能力），合谋价的优势需要惩罚即时（$\delta \to 1$），而惩罚的执行需要对手价格实时可观测（市场透明）。结果：\important{算法合谋不需要协议也不需要教唆，透明、高频、会学三个条件在数字市场同时满足时，超竞争价格是算法交互学习的均衡结果}。
\end{derivation}

执法的困难也随之升级：第 6 章已经说明反垄断法以"协议"为责任基础，算法自主收敛的超竞争价格没有意思联络，无法落入"垄断协议"的构成要件；强化学习机制进一步说明，连"设计合谋算法"的意图都不存在，算法只是被部署去"最大化利润"，合谋是最大化过程的副产品。\mn{算法合谋案例锚点：Calvano 等（2020）双 Q-learning 实验，无沟通自发收敛到超竞争价格；三条件速记"透明、高频、会学"。}中国的算法治理对此的回应是事前备案与分类分级（第 12 章）：让监管者看见算法，在部署环节而不是在均衡结果环节设置约束。

\subsection{算法歧视与算法公平}

算法进入定价、招聘、信贷与推荐决策后，公平问题随之而来：算法没有主观恶意，却可能系统性复制并放大历史数据中的偏见。第 12 章考察了算法治理的制度安排，本节追问其经济学机制：偏见从何而来、为何比人类歧视更难治理、治理的约束是什么。

\textbf{算法歧视}的传导机制有三个环节。第一步是训练数据偏误：历史招聘数据中特定群体样本稀少，或历史信贷数据中已经包含了歧视性的审批结果；第二步是错误相关性学习：模型从有偏数据中习得"群体属性与胜任力（或信用）的伪相关"，把统计关联当作因果依据；第三步是输出强化：模型在部署中依据伪相关筛选候选人、核定利率，被筛选掉的群体在数据中进一步减少，偏误在新一轮训练中被再次放大。三个环节构成一个自我强化的闭环，这也是算法歧视区别于人类歧视的第一重属性：\important{规模性}，一次训练、百万次应用，个体的歧视决策变成了系统的制度参数。第二重属性是\important{不透明性}：深度模型的决策过程不可解释，受害者难以举证歧视的存在与因果，传统反歧视法依赖的"差别对待意图"在算法面前无从认定。

治理方向沿三个层面展开。其一，\textbf{公平性约束}：在训练目标中加入公平性惩罚项，要求模型在关键群体间的预测误差分布接近，以性能损失换取歧视风险的下降；其二，\textbf{可解释性要求}：用归因技术回答"这个拒绝决定由哪些特征驱动"，为个体救济提供证据基础；其三，\textbf{审计制度}：算法备案与第三方审计把部署前的审查变成常规程序（第 12 章），中国的深度合成算法备案累计已突破 5100 款（国家互联网信息办公室，2025），事前审查的覆盖面在迅速扩大。第 6 章算法合谋的教训在此重现：算法的系统性行为需要新的治理工具，沿用"人类行为加人类责任"的传统框架，歧视与合谋都将落在责任体系的空当里。

\section{大模型与平台生态}

\subsection{规模报酬与训练成本}

大模型产业组织分析的起点是成本结构：大模型的能力来自何处、成本如何随规模变化，决定了市场的进入壁垒与竞争格局。本节先给出能力与规模的数量关系，再把它翻译成产业组织的语言。

大模型能力的经验规律是\textbf{scaling law}：模型性能（以损失函数值 $L$ 度量）随参数量 $N$、训练数据量 $D$ 与训练算力 $C$ 的扩大按幂律改善：

\[
L(N) \approx L_0 + \left(\frac{N_0}{N}\right)^{\alpha}, \qquad \alpha \text{ 的经验估计约 } 0.05\text{--}0.1,
\]

幂指数 $\alpha$ 远小于 1 的含义是：性能每改进一个百分点，参数量需要按更大的幅度扩张。训练算力成本与参数量近似成正比（$C_{\text{训练}} \approx k N$），于是性能改进的边际成本随模型规模急剧攀升：

\[
\frac{dC}{d(\text{性能提升})} \propto N^{\alpha+1}.
\]

$\alpha + 1$ 大于 1，同样的性能增量在更大规模上需要付出更高的算力成本，能力竞赛的成本曲线是凸的。这一成本结构的第一层含义是\important{供给方的规模经济极强}：训练固定成本以亿美元计，而推理（使用）的边际成本很低，平均成本随用户量持续下降，大模型行业天然趋向寡头结构（第 3 章）。第二层含义是\textbf{数据要素的再度出场}：模型能力受训练数据质量的约束，高质量数据的边际收益在特定规模后递减，合成数据的替代效应是前沿争论，第 7 章数据要素的分析框架直接适用于大模型时代：厂商的数据采集激励、训练数据授权的版权争议、数据共享的福利分析，都是数据要素市场建设在大模型时代的具体延伸。

成本曲线的形状并不意味着壁垒不可动摇，工程创新持续下移成本曲线，DeepSeek 的迭代提供了直接证据。2024 年底的 V3 以约 557.6 万美元训练出接近顶尖闭源模型的能力（深度求索，2025），压低的是训练端的常数项；2026 年 4 月发布的 DeepSeek-V4 则把竞争焦点移向部署端：约 1.6 万亿参数的混合专家架构（单次推理激活约 490 亿参数）在第三方评测中达到 2025 年版 GPT-5 的综合水平，与当时最前沿模型相差约 8 个月（美国国家标准与技术研究院，2026），而其稀疏注意力设计使百万上下文场景下每 token 的推理算力约为前代的四分之一（深度求索，2026），多数基准上的调用成本低于同等性能的闭源参照模型。需要准确理解这一对比的经济学含义：DeepSeek 没有改变 scaling law 的幂律斜率，能力向 GPT-5 水平靠拢仍然依靠参数规模的万亿级扩张；它改变的是成本函数的位置：架构与工程优化压低常数项，使同等能力可以以更低的算力与价格获得。这意味着规模租金与工程租金并存：既有厂商的算力积累构成壁垒，但算法与工程层面的创新可以绕过部分壁垒，两者共同决定大模型行业的进入条件。这个结论可以落在 scaling law 的参数结构上。

\begin{derivation}{Scaling law 与进入壁垒}{}
\textbf{第一步，从幂律到目标成本。}scaling law 给出
\[
L(N) = L_0 + \left(\frac{N_0}{N}\right)^{\alpha}, \qquad \alpha \in (0.05,\,0.1).
\]
变形并取倒数：要达到目标性能 $L$，所需参数量为
\[
N = N_0\,(L - L_0)^{-1/\alpha}.
\]
由于 $\alpha$ 远小于 1，指数 $1/\alpha$ 远大于 1：性能差距每缩小一点，参数量按幂次放大。\textbf{第二步，把参数量翻译成训练成本。}训练算力与参数量近似成正比，$C = kN$，于是
\[
C = k N_0\,(L - L_0)^{-1/\alpha}.
\]
性能离前沿越近（$L - L_0$ 越小），训练成本按 $(L-L_0)^{-1/\alpha}$ 幂律攀升。设 $\alpha = 0.05$，则 $1/\alpha = 20$：性能差距缩小 1 倍，训练成本扩大 $2^{20}$ 倍，超过一百万倍。\textbf{第三步，进入壁垒的刻画。}新进入者要在目标性能上进入，必须一次性支付固定成本 $C$，而 $C$ 是性能缺口的凸函数：
\[
\frac{dC}{d(\text{性能提升})} \propto N^{\alpha+1} > N.
\]
结果：\important{能力竞赛的成本曲线是凸的，进入前沿的固定成本随性能缺口幂律攀升，供给方规模经济与进入壁垒内生于 scaling law 的参数结构}；工程创新（压低常数项 $k$）与数据渠道（改变有效 $N_0$）可以下移成本曲线，但不改变幂律的斜率，这是 DeepSeek 压低训练成本却不改变行业集中逻辑的原因。
\end{derivation}

中国的算力基础设施为这一竞争提供了底座，2025 年底智能算力规模超过 1500\linebreak EFLOPS，万卡智算集群建成 42 个（工业和信息化部，2025），智能算力的扩张速度居全球前列。\mn{Scaling law：性能进一步、账单成倍加码；训练贵、推理贱，天然趋向寡头。}

\subsection{开源与闭源竞争}

大模型行业的竞争格局中，最核心的分歧不是技术路线，而是开放策略：模型权重是否公开、下游能否自由商用。闭源模型（OpenAI、谷歌）以 API 收费回收训练投资；开源模型（Meta 的 Llama 系列、DeepSeek）以免费发布换取生态。两种策略的经济逻辑需要在统一的产业组织框架中比较。

开源的策略逻辑有三条。其一，\textbf{生态锁定}：开源模型吸引开发者与二次开发，形成以其为基座的工具链生态，基座地位一旦形成，竞争者替换基座的成本极高（第 3 章切换成本）。其二，\textbf{削弱对手}：高质量开源模型压缩闭源模型的溢价空间，这是"商品化你的互补品"策略的直接应用：闭源厂商的模型是开源厂商的下游创新投入品，投入品的价格被打下来，下游创新（应用生态）的租金就留在开源厂商一侧。其三，\textbf{安全与信任}：可审查性获得监管机构与政企客户的信任溢价，敏感行业（金融、政务）的本地化部署天然偏好开源。其中第二条"商品化你的互补品"可以形式化。

\begin{derivation}{开源与闭源的纵向竞争}{}
设上游闭源模型厂商的模型能力为 $v$，模型以价格 $p$ 卖给下游应用开发商；开发者应用的价值为 $v + e$（$e$ 为应用自身贡献，归一化 $e=0$），即开发者的毛利润为 $v - p$。开放源代码模型以能力 $v_0$ 进入（$0 < v_0 < v$），开源模型的价格为零。\textbf{第一步，开发者的参与约束。}开发者选择闭源模型的条件是 $v - p \ge v_0$，闭源厂商的最优定价把约束取等号：
\[
p^* = v - v_0.
\]
\textbf{第二步，闭源厂商的租值。}设参与开发者数量为 $D$，闭源厂商的利润为
\[
\pi_{\text{闭}} = D \cdot (v - v_0).
\]
结果：\important{开源模型把闭源模型的定价上限钉在能力差距 $v - v_0$ 上，闭源租值随开源能力的逼近线性收缩}。反过来看，开源模型是闭源模型的下游互补品，把互补品价格打到零，开发者的剩余 $v - p$ 从 $0$ 提高到 $v_0$，上游租值向下游转移，应用生态（开发者数量与创新密度）随之扩张。这一策略的经济实质是纵向链条上的租值再分配：开源方放弃模型层租金，换取应用层的生态份额；闭源方要维持定价空间，只能维持并扩大能力差距。
\end{derivation}

福利分析框架沿三个维度展开：\textbf{可获得性}，开源提升模型的可及性，消费者剩余上升；\textbf{创新激励}，开源加速下游应用创新，但闭源模式下训练投资的回报保障更强，两类模式各有其创新分工；\textbf{安全}，开源使对齐与安全控制从集中式转为分布式，安全风险的监控成本上升。\mn{开源 vs 闭源的分析框架（复试素材）：福利维度（可获得性、创新激励）、安全维度（对齐控制）、产业维度（生态锁定、市场势力），三条线展开。}产业维度的结论与第 6 章的市场势力分析一致：闭源叠加算力优势的市场集中度更高，开源是竞争格局中的结构性约束力量。

中国的开源生态为这套框架提供了最新注脚：开源模型累计下载量突破 100 亿次、居全球首位（工业和信息化部，2026），2026 年 4 月以 MIT 许可证开源的 DeepSeek-V4 是当时参数规模最大的开源权重模型，发布后迅速成为全球开发者的基础模型之一，闭源模型的 API 定价随之出现多轮下调。开源与闭源在中国市场的竞争不是零和博弈：开源模型压缩基础能力的溢价，闭源模型在安全合规、垂直领域定制上保持差异化空间，两者的边界由下游应用的定制化程度决定。这一分工格局的含义是：大模型竞争不是单一层次的技术竞赛，基础能力层的开源挤压与垂直服务层的闭源差异化并行发生，分析竞争格局必须首先识别竞争的层次。

\subsection{大模型平台的生态竞争与 API 经济}

大模型的价值不在模型本身，而在模型能力的传递方式：绝大多数企业与开发者并不训练模型，而是通过应用程序接口（API）调用模型能力，模型因此成为一种按量计费的标准化服务。\textbf{模型即服务}（Model as a Service，MaaS）的兴起，使大模型行业的竞争从模型性能延伸到生态组织，第 4 章双边市场的框架在此重演。

API 经济的产业结构分为三层：上游模型厂商（提供基础模型与算力）、中游平台服务（模型托管、微调工具、安全合规）、下游应用开发商（面向最终用户的 AI 应用）。这一分层结构的两个特征是分析的入口。其一，\textbf{按调用量计费}：API 以 token 为计量单位，训练成本在模型厂商处一次性沉淀，推理成本随调用量线性发生，定价结构因此介于"一次性授权"与"持续订阅"之间，模型厂商有强烈的动机通过低价扩大调用规模以摊薄训练固定成本。其二，\textbf{供应商与竞争者的双重身份}：模型厂商既向应用开发商供应能力，又可能亲自下场做应用（如官方助手类产品），供应商对下游市场的进入威胁构成应用层的创新约束，这是第 6 章自我优待问题在大模型行业的重演。\mn{API 经济的分析入口：成本侧看 token 计价的规模经济摊薄，结构侧看模型厂商的供应商与竞争者双重身份，两面都指向生态竞争而非单品竞争。}生态竞争的展开方式与第 4 章的鸡生蛋逻辑同构：模型平台以免费额度与低价 API 补贴开发者侧，开发者应用的质量吸引用户侧，用户数据与反馈再反哺模型改进，三者的正反馈循环决定平台的生态规模。中国市场的调用量数据印证了这一循环的运转速度：2025 年底日均 token 调用量超过 100 万亿（国家互联网信息办公室，2025），调用规模的扩张正在把大模型从演示性产品变成基础设施性服务。

API 经济的政策含义落在市场势力与创新激励的平衡上：模型层的高集中度若向下游传导（通过 API 定价、数据回流条款、生态锁定），第 6 章的分析工具即可直接调用；反过来，模型层的开源竞争（14.2.2）又在持续压低 API 价格，产业内生的竞争约束与反垄断规制在此形成互补。API 经济因此是反垄断与创新激励两个主题的交汇案例：同一个市场结构中，规模经济指向集中，生态竞争指向开放，结论取决于对哪一层市场的界定。

\section{Web3 与代币经济}

\subsection{Web3 的概念}

Web3 是数字经济叙事中争议最大的概念之一：支持者视其为下一代互联网，批评者视其为融资话术。剥离技术词汇后，它的经济学内容其实相当清晰：一次所有权结构的重组。本节先给出概念边界，再分析所有权转移的经济含义。

\begin{definition}[Web3]
Web3 是指以区块链为基础、以用户拥有数据与资产为特征的下一代互联网形态：用户以密码学方式拥有资产与身份，平台的收益分配权向用户与贡献者转移。
\end{definition}

其演进脉络：Web1 是"只读"（门户网站分发内容，用户被动浏览），Web2 是"读写"（用户生成内容，但平台拥有数据与收益分配权），Web3 是"读写 + 拥有"。演进的关键变量不是技术形态，而是\textbf{所有权结构}：Web2 的收益集中于平台，用户贡献内容与数据，平台获得广告与佣金收益，贡献者不分享平台的资本增值；Web3 以\textbf{代币}（token）把平台的经济价值分配给用户与贡献者，代币同时承担三重功能：\textbf{支付手段}（平台内的交易媒介）、\textbf{治理投票权}（参与协议参数的表决）、\textbf{价值贮藏}（分享平台增值的载体）。从第 2 章委托代理的视角看，代币化的实质是把用户从"外部贡献者"变成"剩余索取者"，平台与用户的激励结构从对抗转为部分对齐。

代币的法律与经济性质需要辨析：代币是货币还是资产？按货币三职能（价值尺度、流通手段、贮藏手段）对照，绝大多数代币的价格高波动破坏了价值尺度职能，接受范围有限削弱了流通手段职能，代币更接近"数字资产的产权凭证"而非货币。这与第 10 章比特币货币性的辨析结论一致，并可推广：\important{数字人民币是钱（法定货币的数字化），第三方支付是钱包（支付渠道），代币是资产（产权凭证）}。\mn{辨析答题路径：先报货币三职能，再对照代币逐条排除（高波动破坏价值尺度、接受有限削弱流通手段），落点"资产而非货币"；与数字人民币、第三方支付的定位辨析见第 10 章。}中国监管对这一边界的立场是明确的：2017 年起对代币发行融资（ICO）全面禁止，2021 年进一步整治虚拟货币交易与"挖矿"活动，代币的金融属性被严格限制，而区块链技术本身的数据存证、智能合约等非金融应用则持续获得产业政策支持。技术与金融属性的分离监管，是理解中国 Web3 政策走向的基本框架。

\subsection{代币经济的货币数量方程}

代币经济体的核心问题是估值与激励：一个发行代币的项目，币价由什么支撑，增发代币补贴用户是否可行，币值与增长如何权衡。货币理论为这些问题提供了基准分析工具：费雪方程式。本节把它用于代币经济体，推导通胀设计的权衡。

\begin{derivation}{代币估值与通胀设计}{}
\textbf{代币估值与通胀设计。}设代币流通量为 $M$，流通速度为 $V$；代币经济体内服务以法币计价的价格水平为 $P$，真实产出（服务量）为 $Y$；代币的法币价格为 $Q$。代币流通市值为 $QM$，以法币计价的交易额为 $PY$。\textbf{第一步}，写出代币经济体的费雪方程：每单位代币在每期内周转 $V$ 次，代币市值支撑的交易额为
\[
QMV = PY.
\]
\textbf{第二步}，解出代币价格，即代币估值的基本方程：
\[
Q = \frac{PY}{MV}.
\]
代币价格由三个基本量决定：以法币计价的经济活动规模 $PY$（分子）、流通量与流通速度（分母）。\textbf{第三步}，考察增发的比较静态：对 $M$ 求偏导，
\[
\frac{\partial Q}{\partial M} = -\frac{PY}{M^{2}V} < 0.
\]
结果：\important{在 $PY$ 与 $V$ 不变时，代币增发严格压低币价}，增发比例近似等于币价的贬值比例。
\end{derivation}

这个结果有两层政策含义。其一，\textbf{通胀设计的权衡}：代币增发（$M$ 上升）是代币项目补贴冷启动的标准手段，激励早期参与者正是第 4 章鸡生蛋问题的代币化解决方案；但按比较静态结果，增发在 $PY$ 不变时压低币价，激励通胀与币值稳定不可兼得。代币经济体的"货币政策"需要在增长融资与币值承诺之间权衡，这与第 10 章央行数字货币的分析形成对照：数字人民币的发行与回收服务货币政策目标，币值是法定承诺；代币的币值没有任何承诺，全靠发行纪律与经济活动规模支撑。同一个数量方程在两种制度下给出不同结论：央行以币值稳定与实体经济为目标，代币发行者以自身收益最大化为目标。其二，\textbf{投机与基本面}：若代币需求主要来自投机持有而非平台使用，交易媒介职能弱化，流通速度 $V$ 下降且极不稳定，币价由预期而非基本面支撑。数量方程中的 $Y$ 才是长期价值的锚：投机驱动的币价偏离基本面，等价于"货币超发"式的估值虚高，终将向基本面回归。DeFi 项目的兴衰为此提供了反复上演的例证：以增发代币补贴用户的项目，经济系统在补贴退坡与活跃度下降后普遍进入收缩，币值与活跃度同向变动，正是数量方程中 $PY$ 与 $Q$ 正向关系的直观体现。

\subsection{代币激励的委托代理：质押机制}

代币经济的第二个设计问题是激励的执行：去中心化网络没有雇主监督贡献者，谁来保证验证者诚实记账？区块链的答案是把诚实变成验证者的占优策略：质押机制。本节用第 2 章委托代理框架分析其激励逻辑。

\begin{derivation}{质押机制与激励相容条件}{}
\textbf{质押机制与激励相容条件。}设验证者质押 $S$ 个代币作为抵押品；诚实参与获得验证报酬 $w$；不诚实行为（双重支付、伪造交易）可获得额外收益 $B$；不诚实行为被发现的概率为 $p$；被发现后削减（slashing）比例为 $\varphi$ 的质押，即惩罚额为 $\varphi S$。\textbf{第一步}，分别写出两种行为的期望收益。诚实参与：
\[
E[\text{诚实}] = w.
\]
不诚实行为的期望收益为报酬加作恶收益、减去期望惩罚：
\[
E[\text{不诚实}] = w + B - p\,\varphi S.
\]
\textbf{第二步}，写出激励相容条件：不诚实不占优，即 $E[\text{诚实}] \geq E[\text{不诚实}]$，代入得
\[
w \geq w + B - p\,\varphi S
\quad\Longleftrightarrow\quad
p\,\varphi S \geq B.
\]
\textbf{第三步}，解出所需的最小质押额：
\[
S \geq \frac{B}{p\,\varphi}.
\]
结果：\important{质押额必须不低于作恶收益与期望惩罚系数之比 $B/(p\varphi)$}。
\end{derivation}

激励相容条件的比较静态含义清晰：作恶收益 $B$ 越大（可双重花费的金额越高）、发现概率 $p$ 越低（监测能力弱）、削减比例 $\varphi$ 越小（惩罚软弱），所需质押额越高。质押机制的经济本质是\textbf{可强制执行的抵押品}：不当行为的事后后果被制度化为事前成本，验证者的行为约束不依赖任何人的监督意志，而依赖其自身质押资产的安全。这正是第 2 章道德风险问题的解决思路在去中心化场景中的复现：委托代理模型中的激励相容约束需要第三方强制才能成立，质押把强制力内置于契约本身，代码执行取代了法庭执行。

质押机制与代币激励构成代币经济治理的两翼，两者的分工可以从委托代理框架中看出：代币分配解决的是\textbf{激励问题}（贡献者是否愿意持续贡献，即激励相容约束的满足），质押解决的是\textbf{约束问题}（验证者是否可能作恶，即不当行为的成本内置），一个提供胡萝卜、一个提供大棒。以太坊 2022 年 9 月完成从工作量证明（PoW）到权益证明（PoS）的转换，能耗下降超过 99\%（以太坊基金会，2022），验证者以质押代币取代算力竞赛维持账本安全，这是质押机制从理论到主流的标志性实践。与第 10 章的对照同样重要：PoW 的安全性由算力成本购买，攻击成本是算力租金；PoS 的安全性由质押资产购买，攻击成本是质押损失，两者的共性是安全性始终是一个经济均衡而非技术常数，这是理解"区块链安全"的经济学起点。代币经济的估值与治理设计共同受制于监管环境：代币的发行、交易与质押在任何司法辖区都首先是一个合规问题，制度环境决定这一经济实验的边界。这正是下一节政策讨论的入口。

\section{数字经济政策与展望}

\subsection{中国的政策体系}

中国数字经济的政策体系以\important{三个纲领性文件}为骨架。\textbf{《数字中国建设整体布局规划》}（2023）确立"2522"框架（第 1 章）：夯实数字基础设施与数据资源体系两大基础，推进数字技术与经济、政治、文化、社会、生态文明建设五位一体深度融合，强化数字技术创新体系与数字安全屏障两大能力，优化数字化发展国内国际两个环境。\textbf{"数据要素×"三年行动计划}（2024）以数据要素的乘数效应为主线，推动数据在工业制造、现代农业、金融服务、交通运输等 12 个领域的场景化应用，其经济学定位是第 7 章"数据价值化"的行动方案。\textbf{"人工智能+"行动}（2024 年政府工作报告）将 AI 定位为赋能千行百业的通用技术，与"互联网+"形成代际递进。\mn{政策三骨架记忆：数字中国（2522 总框架）→ 数据要素×（数据要素市场化）→ 人工智能+（技术赋能）。三者关系：总框架定方向，数据要素×建设基础制度，人工智能+部署前沿技术。}政策体系的内在逻辑与本书的理论主线同构：先建设要素市场（数据），再推广通用技术（AI），最终实现数字技术与实体经济的深度融合（第 11 章产业数字化）。三个文件的顺序不是随机的：要素市场化改革是技术扩散的前提，通用技术的全面渗透是要素价值化的归宿，理解这一递进关系，是把握政策体系整体脉络的关键。

\subsection{数据要素市场建设的进展与堵点}

数据要素市场建设（第 7 章）的进展呈现"制度先行、交易滞后"的特征。制度层，三权分置确权框架、数据资产入表、数据交易所体系已经就位；交易层，全国数据交易机构超过 50 家（截至 2025 年 7 月），但场内交易占数据流通总量的比例仅约两成，场外点对点交易仍是主体，多数交易机构的商业化能力薄弱。头部机构的分化印证了"制度先行、交易滞后"的判断：深圳数据交易所累计交易规模突破 150 亿元（截至 2025 年 6 月），已跑通规模化路径，而多数中小交易所尚在探索商业模式。国家数据局 2025 年工作要点提出年内场内交易规模突破 5000 亿元的目标，并明确对交易场所实施统筹优化、严控数量，制度建设的重心开始从"建机构"转向"促交易"。

堵点仍在第 7 章识别的三者：确权成本、定价困难与合规负担。确权环节，数据来源的多元性使三权分置的落地需要逐案界定；定价环节，数据价值高度场景依赖，缺乏可参照的市场基准；合规环节，跨主体数据融合的授权链条冗长。突破方向沿三条线推进：公共数据的授权运营（以公共数据撬动市场供给，公共数据占数据总量的比重高而开发率低）；行业数据空间（行业内的可信共享基础设施，降低同行业企业的互信成本）；数据资产融资的制度配套（质押、证券化，把数据资产的账面价值转化为流动性）。判断进展的理论标尺是第 7 章的效率条件：数据流通的边际成本是否实质性下降，制度建设的成效最终以交易成本衡量，而非以机构数量或交易金额的绝对规模衡量。

\subsection{数字经济与中国式现代化}

人工智能、大模型与代币经济的前沿命题，最终都汇入一个整体问题：数字经济的制度建设如何融入中国式现代化的进程。党的二十大报告把"加快发展数字经济，促进数字经济和实体经济深度融合"列为现代化产业体系建设的重要方向，数字经济因此不只是产业命题，而是现代化道路的制度命题。数字经济的理论主线与中国式现代化的特征在三个接口上对应。\textbf{高质量发展的现代化}：数据要素市场化与数字产业化提供增长新动能（第 8 章增长核算的构成项），2025 年人工智能核心产业规模超过 1.2 万亿元（工业和信息化部，2025），数字产业已是产业体系中最活跃的部门。\textbf{人口规模巨大的现代化}：超大规模市场与网络外部性的结合是中国平台经济的基本盘（第 3 章），数字基础设施的普惠化使服务触达成本在 14 亿人口规模上被摊薄。\textbf{共同富裕的现代化}：数字普惠金融扩大金融服务的覆盖（第 10 章），数字技能的再分配政策缓解技能溢价上升带来的收入分化（第 9 章），数字鸿沟的治理是这一接口的关键约束。三个接口对应的政策文件正是 14.4.1 的三骨架：要素市场对应增长动能，基础设施对应规模优势，普惠导向对应分配公平。

面向未来，本书的分析框架对\important{三个未决问题}保持开放。\textbf{第一，规模报酬递增与均衡的存在性}：数据与网络效应使规模报酬递增成为常态，传统竞争均衡的理论基础（凸性假设）动摇，数字经济需要新的均衡概念与福利标准（第 3、6 章的开放问题）。\textbf{第二，算法经济中的责任与激励}：当配置决策由算法完成，传统的"人加责任"框架失效，算法审计、可解释性与制度设计的边界仍在探索（第 6、12 章）。\textbf{第三，数字红利的全球分配}：数字鸿沟在国家间重演，全球数字治理的规则供给远落后于数字经济的发展速度，产业链安全（第 11 章）与数据跨境流动（第 12 章）的国内制度安排，都要在尚不健全的国际规则环境中运行。这三个问题都指向同一个方向：\important{数字经济的理论建设与制度建设是同一项工程的两面}，经济学的任务不是事后解释数字经济的成功，而是为它的下一个阶段提供制度想象力。

回答这些问题，最终需要把全书的工具装配成一张可调用的地图：一个数字企业从冷启动到垄断的完整旅程，正是全部章节工具的运转现场。第 15 章以数字企业的生命周期为主线，完成这次装配。

\begin{chapterbox}
\textbf{本章考点清单}

\begin{enumerate}
\item AI 的 GPT 判定三标准与生产率时滞 \important{【专硕·核心】}
\item 任务模型的 AI 扩展：认知任务切入、任务分析单位与边界外移就业冲击 $g_N - dx^*/dt$ 的推导 \important{【专硕·核心】}
\item 算法合谋的强化学习机制与三个达成条件（重复博弈 $\delta^*$、时滞 $\delta=e^{-r\Delta t}$、Q-learning 收敛） \important{【专硕·扩展】}
\item 算法歧视的传导闭环与治理方向 \important{【专硕·扩展】}
\item Scaling law 与训练成本的幂律攀升、进入壁垒 $C = kN_0(L-L_0)^{-1/\alpha}$ 的推导、DeepSeek 案例 \important{【专硕·扩展】}
\item 开源闭源竞争的福利框架、纵向竞争模型 $p^* = v - v_0$ 与 API 经济 \important{【专硕·扩展】}
\item Web3 的所有权结构变革与代币三重功能 \important{【专硕·扩展】}
\item 代币数量方程 $Q = PY/(MV)$ 与通胀设计权衡 \important{【专硕·扩展】}
\item 质押机制的激励相容条件 $S \geq B/(p\varphi)$ \important{【专硕·扩展】}
\item 中国政策三骨架（2522、数据要素×、人工智能+）\important{【专硕·核心】}
\item 数据要素市场建设的堵点与标尺 \important{【专硕·核心】}
\item 数字经济与中国式现代化（论述题收束）\important{【专硕·核心】}
\end{enumerate}
\end{chapterbox}
